\section{What counting and structural identities do not establish}
\begin{lemma}[Counting local descriptions \claimid{L01}]
\label{lem:count}
Fix \(n,t\) and an integer \(s\geq0\). At most \(2^{s+1}-1\)
truth tables on \(n\) input bits have \(\loc_t(T)\leq s\).
For a uniformly random such table,
\[
 \Pr[\loc_t(T)\leq s]\leq\min\{1,2^{s+1-2^n}\}.
\]
\end{lemma}
\begin{proof}
There are \(\sum_{j=0}^s2^j=2^{s+1}-1\) descriptions of length at
most \(s\). A deterministic description satisfying the time bound
on every query determines at most one table. Divide the bound by
the total \(2^{2^n}\) possible tables.
\end{proof}
This statement identifies most tables as hard, not SAT's particular
table. Proposition~\ref{prop:short} explains why raw information content
cannot identify the desired computational obstruction. Shannon entropy
also requires a specified distribution; the size of one table is not
its computational hardness.

\begin{proposition}[Self-reduction requires boundary conditions \claimid{P02}]
\label{prop:boundary}
Let \(P\) be Boolean-valued on formulas, allowing Boolean constants.
Suppose for every variable \(z\) occurring in \(F\),
\[
 P(F)=P(F[z=0])\lor P(F[z=1]).
\]
If \(P\) equals direct evaluation on variable-free formulas, then
\(P(F)=\SAT(F)\) for every \(F\). Without that boundary condition,
both constant functions satisfy the recurrence.
\end{proposition}
\begin{proof}
Induct on the number of distinct variables. The boundary condition
proves the base case. Substitution removes \(z\), so induction
identifies both terms on the right as satisfiability answers; their
disjunction is satisfiability of \(F\). For a constant function with
value \(c\), the recurrence is \(c=c\lor c\).
\end{proof}
Uniqueness under identities and boundary conditions does not bound
circuit size. Substituted formulas also have different encoding
lengths: a circuit-family argument must manage these lengths or
supply explicit satisfiability-preserving padding.
Counting identities as independent gates is unjustified because
circuits can share subcomputations.

A deterministic halting computation using space \(s(n)\geq\log n\)
has exponentially many possible configurations in \(s(n)\), for a
fixed standard machine. Counting them gives an exponential
\emph{upper} bound on time. It does not force a computation to visit
all configurations and is not a time lower bound.
